% ============================================================
%  Helicase LXR MIDI Tables
%  The old LXR CC list is gone: parameters are addressed by NRPN =
%  canonical 9-bit target ID.  Maintain the ID ranges here.
%
%  Environments (from preamble.tex): miditableCC, miditableNRPN
%  Row format:   ID range & Meaning & Channel \\
% ============================================================

\subsection*{NRPN address space}

\begin{miditableNRPN}
0 -- 383   & Voice parameters: slot $\times$ 64 + parameter index (slots 0 to 5 = voices 1 to 6) & G or own voice channel \\
384 -- 389 & Voice Morph amounts \texttt{1vm} to \texttt{6vm} & G or own voice channel \\
390        & Retired (the former global sample-rate control); ignored & --- \\
391        & Track 7 decay (\texttt{7dc}), the open hi-hat decay & G or voice 6 channel \\
392 -- 397 & Voice outputs \texttt{1ou} to \texttt{6ou} & G or own voice channel \\
398 -- 403 & Voice effect sends \texttt{1fx} to \texttt{6fx} & G or own voice channel \\
404        & Effect Morph \texttt{fxm} & G only \\
405 -- 411 & Track length, tracks 1 to 7 & G or own voice channel \\
412 -- 418 & Track step scale, tracks 1 to 7 & G or own voice channel \\
419 -- 425 & Track shuffle, tracks 1 to 7 & G or own voice channel \\
426 -- 447 & Unused (ignored) & --- \\
448 -- 510 & Effect parameters: Effect-local index 0 to 62 & G only \\
511        & Off (ignored) & --- \\
\end{miditableNRPN}

\noindent\verify{The individual IDs from 391 to 403 are derived from the documented order
(voice Morph, \texttt{7dc}, outputs, sends) and the documented position of \texttt{fxm} at 404 and the
retired global control at 390. Confirm them against the project's published NRPN chart.}

\vspace{1ex}
\noindent The voice parameter chart (index within each Instrument type) is maintained separately by
the project. The Alkaline voice has 30 parameters (indices 0 to 29); indices 30 to 63 of its block
are unused. For the other types the indices follow the descriptor order of each Instrument.
\verify{Include the per-type voice NRPN chart in the final manual (published separately as the project's NRPN table).}

\subsection*{Parameter value notes}

\begin{miditableCC}
\texttt{osc2\_mod\_type} (Drum \texttt{mod}) & 0 \texttt{fm}, 1 \texttt{mix}, 2 \texttt{snc}, 3 to 127 \texttt{rng} & --- \\
\texttt{osc1\_type} (\texttt{typ})           & 0 \texttt{wav}, 1 \texttt{smp}, 2 \texttt{slp}, 3 \texttt{wtb}, 4 \texttt{pwm} & --- \\
\texttt{osc1\_run} (\texttt{run})            & 0 \texttt{rst}, 1 \texttt{fre}, 2 \texttt{itp}; above 2 acts as \texttt{itp} & --- \\
\texttt{voice\_chain} (\texttt{chn})         & 0 \texttt{T+D}, 1 \texttt{WFd}; above 1 acts as \texttt{WFd} & --- \\
Alkaline waveform                            & 0 \texttt{Saw}, 1 \texttt{Sqr}, 2 \texttt{Tri} & --- \\
Alkaline slide                               & 0 off, 1 on & --- \\
Alkaline tuning                              & 0 to 127, 63 is centre ($\pm$700 cents) & --- \\
Morph amounts                                & 0 to 126 are doubled; 127 gives 255 & --- \\
\end{miditableCC}

\subsection*{Default note map}

For a card with no per-track channel or note settings, the global channel stripes the notes across
the seven tracks at the default pitch, so notes 0 to 6 trigger tracks 1 to 7 in turn and the
pattern repeats every seven notes. Set a track's root note on the STEP page (\param{nte}) to give
it a chromatic range.

\verify{The earlier LXR responded on notes 36 to 42 by default. Confirm the Helicase default.}
